\documentclass[aspectratio=169,11pt]{beamer}

% Minimal theme with clean look
\usetheme{default}
\usecolortheme{seahorse}
\setbeamertemplate{navigation symbols}{}
\setbeamertemplate{footline}[frame number]
\setbeamertemplate{frametitle}[default][center]
\setbeamerfont{frametitle}{size=\Large,series=\bfseries}
\setbeamercolor{frametitle}{fg=black,bg=white}

% Remove shadows and reduce visual clutter
\setbeamertemplate{blocks}[default]
\setbeamertemplate{itemize items}[circle]

% Packages
\usepackage{amsmath}
\usepackage{amssymb}
\usepackage{booktabs}
\usepackage{graphicx}
\usepackage{xcolor}
\usepackage{pifont}

% Colors
\definecolor{darkblue}{RGB}{0,51,102}
\definecolor{medblue}{RGB}{51,102,153}
\definecolor{lightblue}{RGB}{153,204,255}
\definecolor{darkred}{RGB}{153,0,0}
\definecolor{darkgreen}{RGB}{0,102,51}

% Custom commands
\newcommand{\good}{\textcolor{darkgreen}{\ding{51}}}
\newcommand{\bad}{\textcolor{darkred}{\ding{55}}}
\newcommand{\maybe}{\textcolor{orange}{!}}

\title{Quantum-Inspired Metrics for LLM Distribution Testing:\\
Initial Results and Discoveries}
\author{Model Equality Testing Research}
\date{\today}

\begin{document}

% Title slide
\begin{frame}
\titlepage
\end{frame}

% Motivation
\begin{frame}{Motivation: Beyond Binary Detection}

\textbf{Existing capability:} Detect distribution changes with MMD
\begin{itemize}
  \item Character-level kernel (Hamming)
  \item Statistical power, rigorous p-values
  \item Answer: \emph{"Did the distribution change?"} \good
\end{itemize}

\vspace{1em}

\textbf{Goal:} Interpretability — understand \emph{how} distributions differ
\begin{itemize}
  \item Quantum-inspired metrics (trace distance, Von Neumann entropy)
  \item Semantic embeddings to capture meaning
  \item Answer: \emph{"What changed semantically?"}
\end{itemize}

\end{frame}

% Quantum metrics approach
\begin{frame}{Quantum-Inspired Approach}

\textbf{Core idea:} Represent distribution of LLM outputs as quantum state

\vspace{0.5em}

\textbf{From completions to density matrices:}
\begin{enumerate}
  \item Start with N completions from model (text strings)
  \item Embed each completion: $\mathbf{x}_i \in \mathbb{R}^{768}$ (using MPNet or EmbeddingGemma)
  \item Normalize: $\psi_i = \mathbf{x}_i / ||\mathbf{x}_i||$ (unit vectors on sphere)
  \item Construct Gram matrix: $G_{ij} = \langle \psi_i | \psi_j \rangle$ (similarity between samples)
  \item Density matrix: $\rho = \frac{1}{N} G$ (normalized to unit trace)
\end{enumerate}

\vspace{0.5em}

\textbf{What the density matrix captures:}
\begin{itemize}
  \item Diagonal: individual sample "energies"
  \item Off-diagonal: correlations/overlaps between samples
  \item Eigenvalues: principal semantic dimensions
  \item Trace = 1 (normalization), Hermitian, positive semi-definite
\end{itemize}

\end{frame}

% Why quantum
\begin{frame}{Why Quantum Formalism?}

\textbf{Classical probability} (single distribution):
\begin{itemize}
  \item Represented by probability vector: $\mathbf{p} = (p_1, p_2, \ldots, p_k)$
  \item Events are mutually exclusive (one outcome at a time)
  \item No notion of "superposition" or "interference"
\end{itemize}

\vspace{0.5em}

\textbf{Quantum probability} (density matrix formalism):
\begin{itemize}
  \item Represented by density matrix: $\rho \in \mathbb{C}^{d \times d}$
  \item Captures correlations and "coherence" between states
  \item Natural geometric structure (Hilbert space)
  \item Rich set of information-theoretic metrics
\end{itemize}

\vspace{0.5em}

\textbf{Analogy:} Classical treats samples as points; quantum treats them as overlapping waves

\vspace{0.5em}

\textbf{Hypothesis:} Semantic embeddings have wave-like structure (overlaps matter)
\begin{itemize}
  \item Quantum metrics might capture subtle correlations
  \item Classical p-values miss geometric relationships
\end{itemize}

\end{frame}

% What the metrics mean
\begin{frame}{What Do the Quantum Metrics Mean?}

\textbf{1. Trace Distance} $D(\rho_A, \rho_B) = \frac{1}{2}\text{Tr}|\rho_A - \rho_B|$
\begin{itemize}
  \item \textbf{Range:} [0, 1]
  \item \textbf{Intuition:} Maximum probability of distinguishing two distributions
  \item 0 = identical distributions (indistinguishable)
  \item 1 = orthogonal distributions (perfectly separable)
  \item \textbf{Geometric:} Distance in the space of quantum states
\end{itemize}

\vspace{0.5em}

\textbf{2. Von Neumann Entropy} $S(\rho) = -\text{Tr}(\rho \log \rho)$
\begin{itemize}
  \item \textbf{Range:} [0, $\log d$] where $d$ is embedding dimension
  \item \textbf{Intuition:} Diversity/uncertainty of the distribution
  \item Low entropy = outputs cluster tightly (deterministic)
  \item High entropy = outputs spread broadly (diverse)
  \item \textbf{Interpretability:} Does quantization reduce diversity?
\end{itemize}

\vspace{0.5em}

\textbf{3. Quantum Relative Entropy} $D(\rho_A || \rho_B) = \text{Tr}(\rho_A \log \rho_A - \rho_A \log \rho_B)$
\begin{itemize}
  \item \textbf{Range:} [0, $\infty$)
  \item \textbf{Intuition:} Information divergence (quantum KL divergence)
  \item Measures information gained when using true dist. $\rho_A$ vs approximate $\rho_B$
  \item \textbf{Asymmetric:} $D(\rho_A || \rho_B) \neq D(\rho_B || \rho_A)$
  \item \textbf{Interpretability:} Information content shift from implementation change
\end{itemize}

\end{frame}

% Why we thought it would work
\begin{frame}{Why We Expected This to Work}

\textbf{Hypothesis:} Quantum metrics capture semantic structure classical tests miss

\vspace{0.5em}

\textbf{Expected advantages:}

\begin{enumerate}
  \item \textbf{Geometric interpretability}
  \begin{itemize}
    \item Trace distance: magnitude of shift
    \item Entropy: diversity changes
    \item Relative entropy: information gained/lost
  \end{itemize}

  \vspace{0.3em}

  \item \textbf{Sample efficiency}
  \begin{itemize}
    \item Density matrix captures correlations
    \item Might detect subtle shifts with fewer samples
    \item Initial results seemed promising (0.32 distance at n=100)
  \end{itemize}

  \vspace{0.3em}

  \item \textbf{Characterization beyond detection}
  \begin{itemize}
    \item Not just "distributions differ" but \emph{how} they differ
    \item Eigenvalue spectrum reveals structure
    \item Direction of change in semantic space
  \end{itemize}
\end{enumerate}

\vspace{0.5em}

\textbf{Reality check:} Sanity tests revealed fundamental issues...

\end{frame}

% Initial results
\begin{frame}{Initial Results: Promising but Puzzling}

\textbf{Test case:} Llama-3-8B fp32 vs int8 quantization

\vspace{0.5em}

\begin{center}
\begin{tabular}{lcc}
\toprule
\textbf{Metric} & \textbf{n=100} & \textbf{n=1000} \\
\midrule
Trace distance & 0.321 & 0.379 \\
Von Neumann divergence & 0.081 & 0.013 \\
QRE (symmetric) & 0.003 & 0.001 \\
\midrule
MMD (Hamming) p-value & 0.28 & \textbf{0.01} \good \\
\bottomrule
\end{tabular}
\end{center}

\vspace{0.5em}

\textbf{Initial interpretation:} Trace distance detects difference (0.32--0.38) where MMD lacks power at n=100

\vspace{0.5em}

\textbf{Problem:} Is trace distance more sample-efficient?

\end{frame}

% Sanity check failure
\begin{frame}{Critical Finding: Sanity Check Failure}

\textbf{Sanity test:} fp32 vs fp32 (same distribution, two samples)

\vspace{0.5em}

\begin{center}
\begin{tabular}{lccc}
\toprule
\textbf{Metric} & \textbf{fp32 vs int8} & \textbf{fp32 vs fp32} & \textbf{Status} \\
\midrule
Trace distance (n=100) & 0.321 & \textbf{0.332} & \bad \\
Trace distance (n=1000) & 0.379 & \textbf{0.371} & \bad \\
MMD p-value (n=100) & 0.28 & \textbf{0.81} & \good \\
MMD p-value (n=1000) & 0.01 & \textbf{0.38} & \good \\
\bottomrule
\end{tabular}
\end{center}

\vspace{1em}

\textbf{Conclusion:} Trace distance shows \emph{no discrimination} — identical values for different and same distributions!

\vspace{0.5em}

\textbf{Root cause:} Sample-size artifacts in Gram matrix construction

\end{frame}

% Sample size artifact
\begin{frame}{Root Cause: Sample-Size Dependence}

\textbf{Validation on synthetic data} (two samples from N(0,I)):

\vspace{0.5em}

\begin{center}
\begin{tabular}{cc}
\toprule
\textbf{Sample Size} & \textbf{Trace Distance} \\
\midrule
50 & 0.146 \\
100 & 0.212 \\
200 & 0.303 \\
500 & 0.467 \\
1000 & \textbf{0.636} \\
\bottomrule
\end{tabular}
\end{center}

\vspace{0.5em}

\textbf{Expected:} Distance $\to$ 0 as n increases (better estimates)

\textbf{Observed:} Distance $\to$ 1 as n increases \bad

\vspace{0.5em}

\textbf{Issue:} Gram matrix size (N $\times$ N) grows with sample size, creating systematic bias

\end{frame}

% PCA fix
\begin{frame}{Attempted Fix: Dimensionality Reduction}

\textbf{Approach:} Project to fixed k dimensions before constructing density matrix

\vspace{0.5em}

\textbf{Synthetic data (N(0,I) vs N(0,I)):}
\begin{itemize}
  \item Original: 0.209 distance \bad
  \item PCA k=50: \textbf{0.004 distance} \good
  \item Sample-size artifact eliminated! \good
\end{itemize}

\vspace{0.5em}

\textbf{Real LLM data (fp32 vs int8, n=100):}
\begin{itemize}
  \item fp32 vs fp32: 0.076
  \item fp32 vs int8: 0.081
  \item Discrimination: 7\% effect size
\end{itemize}

\vspace{0.5em}

\textbf{Status:} Math fixed, but discrimination too weak for practical use

\textbf{Verdict:} Embedding space lacks signal for quantization effects

\end{frame}

% Pivot to interpretability
\begin{frame}{Pivot: Semantic Embeddings for Interpretability}

\textbf{Reframing the question:}

\vspace{0.5em}

Why do semantic embeddings fail to separate fp32 vs int8?

$\Rightarrow$ Maybe \emph{semantics are preserved} despite token-level changes!

\vspace{1em}

\textbf{Two-level testing framework:}

\begin{enumerate}
  \item \textbf{Token-level (MMD Hamming):} Detect implementation changes
  \begin{itemize}
    \item Sensitive to exact tokens, word choice, phrasing
    \item Answer: "Did the distribution change?" (p-value)
  \end{itemize}

  \vspace{0.5em}

  \item \textbf{Semantic-level (EmbeddingGemma + t-SNE):} Assess meaningful impact
  \begin{itemize}
    \item Preserves meaning, discards surface form
    \item Answer: "Does the change matter semantically?" (separation ratio)
  \end{itemize}
\end{enumerate}

\vspace{0.5em}

\textbf{Hypothesis:} Implementation changes can be statistically detectable without being semantically meaningful

\end{frame}

% Semantic invariance results
\begin{frame}{Discovery: Semantic Invariance Under Implementation Changes}

\textbf{Llama-3-8B vs Mistral-7B comparison} (Wikipedia continuations, n=100):

\vspace{0.5em}

\begin{center}
\begin{tabular}{lcc}
\toprule
\textbf{Comparison} & \textbf{Token-level (MMD)} & \textbf{Semantic (t-SNE)} \\
\midrule
fp32 vs int8 (quant) & p=0.01 (detected) & 0.0 separation \\
fp32 vs fp32 (sanity) & p=0.38 (correct) & 0.0 separation \\
Llama vs Mistral (arch) & (not tested) & \textbf{<0.4 separation} \\
\bottomrule
\end{tabular}
\end{center}

\vspace{0.5em}

\textbf{Separation ratio:} (center distance) / (cluster spread), threshold 2.0 for meaningful separation

\vspace{1em}

\textbf{Key finding:} Quantization changes token distributions but preserves semantic content

\textbf{Interpretation:} Both MMD and embeddings are correct — they measure different abstraction levels

\end{frame}

% Task diversity
\begin{frame}{Task Diversity: Pattern Holds Across Domains}

\textbf{Llama-3-8B vs Mistral-7B across task types:}

\vspace{0.5em}

\begin{center}
\begin{tabular}{llc}
\toprule
\textbf{Dataset} & \textbf{Task Type} & \textbf{Separation Ratio} \\
\midrule
Wikipedia (en) & Factual continuation & 0.40 \\
HumanEval & Code generation & 0.26 \\
UltraChat & Conversational & 0.22 \\
\bottomrule
\end{tabular}
\end{center}

\vspace{1em}

\textbf{Counterintuitive:} More open-ended tasks show \emph{stronger} convergence

\vspace{0.5em}

\textbf{Interpretation:} Different architectures converge to similar semantic strategies regardless of task constraint

\vspace{0.5em}

All ratios $<$ 2.0 threshold $\Rightarrow$ no meaningful semantic separation

\end{frame}

% Cross-language analysis
\begin{frame}{Cross-Language Analysis: Training Bias Revealed}

\textbf{Hypothesis:} Convergence driven by training data abundance?

\vspace{0.5em}

\textbf{Test:} Llama-3-8B (Meta/US) vs Mistral-7B (Mistral AI/France) across 5 languages

\vspace{0.5em}

\begin{center}
\begin{tabular}{lcc}
\toprule
\textbf{Language} & \textbf{Separation Ratio} & \textbf{Interpretation} \\
\midrule
\textbf{French} & \textbf{0.19} & \textcolor{darkgreen}{Strongest convergence} \\
English & 0.40 & Moderate \\
Russian & 0.40 & Moderate \\
German & 0.71 & Some divergence \\
\textbf{Spanish} & \textbf{1.00} & \textcolor{darkred}{Highest divergence} \\
\bottomrule
\end{tabular}
\end{center}

\vspace{0.5em}

\textbf{Pattern:} 5$\times$ variation across languages (0.19 $\to$ 1.00)

Language variation $>$ task variation in English!

\end{frame}

% Verbosity hypothesis
\begin{frame}{Alternative Hypothesis: Linguistic Verbosity?}

\textbf{Question:} Is French convergence due to French being more verbose/constrained?

\vspace{0.5em}

\begin{center}
\begin{tabular}{lccc}
\toprule
\textbf{Language} & \textbf{Avg Length} & \textbf{Verbosity Rank} & \textbf{Separation} \\
\midrule
English & 246.7 chars & 1 (most) & 0.40 \\
Spanish & 172.2 chars & 2 & \textbf{1.00} \\
\textbf{French} & \textbf{170.8 chars} & \textbf{3} & \textbf{0.19} \\
German & 170.0 chars & 4 & 0.71 \\
Russian & 146.9 chars & 5 (least) & 0.40 \\
\bottomrule
\end{tabular}
\end{center}

\vspace{0.5em}

\textbf{Correlation:} -0.141 (essentially none)

\vspace{0.5em}

\textbf{Spanish vs French:} Nearly identical verbosity (172 vs 171 chars), opposite convergence patterns (1.00 vs 0.19)

\vspace{0.5em}

\textbf{Verdict:} \bad\ Verbosity does not predict convergence

\end{frame}

% Training data convergence
\begin{frame}{The Spanish Paradox: Length $\neq$ Semantic Convergence}

\textbf{Per-model analysis reveals both invested in French:}

\vspace{0.5em}

\begin{center}
\begin{tabular}{lccc}
\toprule
\textbf{Language} & \textbf{Llama std} & \textbf{Mistral std} & \textbf{Separation} \\
\midrule
\textbf{French} & \textbf{14.3} (lowest!) & 21.0 & 0.19 \\
Spanish & 13.5 & 21.6 & \textbf{1.00} \\
Russian & 21.3 & 17.7 & 0.40 \\
English & 50.7 & 33.0 & 0.40 \\
German & 23.2 & 32.7 & 0.71 \\
\bottomrule
\end{tabular}
\end{center}

\vspace{0.5em}

\textbf{Key observations:}
\begin{itemize}
  \item Llama's French std (14.3) is lowest across \emph{all} languages — even lower than English!
  \item Mistral's French std 20\% lower than other languages
  \item Both models show French optimization \good
\end{itemize}

\vspace{0.5em}

\textbf{The paradox:} Spanish and French have similar std (13.5 vs 14.3) but opposite convergence (1.00 vs 0.19)

\end{frame}

% Refined hypothesis
\begin{frame}{Refined Hypothesis: Training Data Convergence}

\textbf{Why French converges (0.19):}
\begin{itemize}
  \item Both companies used similar \emph{canonical sources}
  \item French Wikipedia, standardized literature
  \item Parisian French dominates formal writing (homogeneous)
  \item Strategic importance: 300M speakers, EU operations
\end{itemize}

\vspace{0.5em}

\textbf{Why Spanish diverges (1.00):}
\begin{itemize}
  \item \emph{Dialectal variation}: Spain vs Mexico vs Argentina vs Colombia
  \item Models made different regional/dialectal choices
  \item No single dominant variant in training corpora
  \item Produce similar-length text (172 chars) but different semantic styles
\end{itemize}

\vspace{0.5em}

\textbf{Revised framework:}

High convergence $\Rightarrow$ models used similar canonical sources

High divergence $\Rightarrow$ models made different corpus/dialect choices

\vspace{0.5em}

\textbf{Pattern:} Separation ratios reveal \emph{training corpus standardization}, not just organizational bias

\end{frame}

% Decision framework
\begin{frame}{Practical Framework: Two-Tier Distribution Testing}

\begin{center}
\begin{tabular}{cccl}
\toprule
\textbf{Statistical} & \textbf{Semantic} & \textbf{Interpretation} & \textbf{Action} \\
\midrule
No change & — & Identical distributions & \good\ Safe \\
Changed & Preserved & Implementation shift & \good\ Likely safe \\
Changed & Changed & Meaningful shift & \maybe\ Investigate \\
\bottomrule
\end{tabular}
\end{center}

\vspace{1em}

\textbf{Example application:}
\begin{itemize}
  \item \textbf{Quantization (fp32 $\to$ int8):} Statistical change + semantic preservation
  \begin{itemize}
    \item MMD p=0.01 (detected) + separation 0.0 (preserved)
    \item \good\ Deploy with confidence
  \end{itemize}

  \vspace{0.5em}

  \item \textbf{Model swap (Llama $\to$ Mistral):}
  \begin{itemize}
    \item French: separation 0.19 $\Rightarrow$ \good\ Interchangeable
    \item Spanish: separation 1.00 $\Rightarrow$ \maybe\ Test which dialect matches use case
  \end{itemize}
\end{itemize}

\end{frame}

% Key findings summary
\begin{frame}{Summary: Key Findings}

\textbf{1. Quantum metrics for detection:} \bad\ Not viable
\begin{itemize}
  \item Trace distance fails sanity checks (sample-size artifacts)
  \item PCA fix weakens discrimination (7\% effect)
  \item Classical MMD superior for statistical testing
\end{itemize}

\vspace{0.5em}

\textbf{2. Semantic embeddings reveal invariance:} \good\ Valuable insight
\begin{itemize}
  \item Implementation changes preserve semantic content
  \item Quantization (fp32 vs int8): detected but semantics preserved
  \item Architecture (Llama vs Mistral): semantic convergence across tasks
\end{itemize}

\vspace{0.5em}

\textbf{3. Language-dependent convergence:} \good\ Training corpus insight
\begin{itemize}
  \item 5$\times$ variation across languages (French 0.19 $\to$ Spanish 1.00)
  \item Both models optimized French using similar canonical sources
  \item Spanish divergence driven by dialectal variation in training
  \item Separation ratios reveal corpus standardization
\end{itemize}

\end{frame}

% Methodological contributions
\begin{frame}{Methodological Contributions}

\textbf{1. Two-level testing framework validated:}
\begin{itemize}
  \item Token-level (MMD): Statistical detection with p-values
  \item Semantic-level (embeddings): Impact assessment with separation ratios
  \item Provides both "did it change?" and "does it matter?" answers
\end{itemize}

\vspace{0.5em}

\textbf{2. Separation ratios as training corpus diagnostics:}
\begin{itemize}
  \item Low ratios $\Rightarrow$ similar canonical sources (French 0.19)
  \item High ratios $\Rightarrow$ different corpus/dialect choices (Spanish 1.00)
  \item Cross-language profiles reveal organizational priorities
\end{itemize}

\vspace{0.5em}

\textbf{3. Negative results have value:}
\begin{itemize}
  \item Quantum metrics "failure" revealed semantic preservation
  \item Journey to understand why they didn't discriminate led to insights
  \item Both MMD and embeddings correct — different abstraction levels
\end{itemize}

\end{frame}

% Future work
\begin{frame}{Future Directions}

\textbf{1. Validate organizational bias hypothesis:}
\begin{itemize}
  \item Compare two French models (expect low French separation)
  \item Compare two US models (expect low English separation)
  \item Compare Chinese models on Chinese vs English
\end{itemize}

\vspace{0.5em}

\textbf{2. Dialectal variation analysis:}
\begin{itemize}
  \item Identify which Spanish variant each model prefers
  \item Test Arabic (MSA vs Egyptian vs Gulf)
  \item Portuguese (Brazil vs Portugal)
  \item Chinese (Simplified vs Traditional)
\end{itemize}

\vspace{0.5em}

\textbf{3. Extend to other distribution shifts:}
\begin{itemize}
  \item Watermarking: semantic preservation?
  \item Fine-tuning: which dimensions change?
  \item Prompt sensitivity: semantic stability?
\end{itemize}

\vspace{0.5em}

\textbf{4. Decision framework validation:}
\begin{itemize}
  \item Establish separation ratio thresholds
  \item User studies: does semantic preservation predict user perception?
\end{itemize}

\end{frame}

% Conclusion
\begin{frame}{Conclusions}

\textbf{Original goal:} Use quantum metrics for interpretable LLM distribution characterization

\vspace{0.5em}

\textbf{What we learned:}

\begin{enumerate}
  \item Quantum metrics don't improve on classical methods for \emph{detection}
  \begin{itemize}
    \item Sample-size artifacts difficult to eliminate
    \item Embedding space issues fundamental
  \end{itemize}

  \vspace{0.5em}

  \item Semantic embeddings reveal \emph{what doesn't change}
  \begin{itemize}
    \item Implementation differences preserve semantic content
    \item Two-level framework: detection + impact assessment
  \end{itemize}

  \vspace{0.5em}

  \item Cross-language analysis reveals training strategies
  \begin{itemize}
    \item Convergence reflects corpus standardization
    \item Divergence reflects dialectal choices
    \item Separation ratios diagnostic tool
  \end{itemize}
\end{enumerate}

\vspace{0.5em}

\textbf{Practical value:} Framework for deployment decisions (quantization safety, model selection)

\end{frame}

\end{document}
